\subsection{Normalen- und Koordinatenform einer Ebene}
\label{subsec:analytic_geometry_plane_normal_form}

Statt zwei Bewegungsrichtungen kann eine Ebene durch eine einzige senkrechte
Richtung charakterisiert werden. Diese Beschreibung bereitet insbesondere die
Abstandsformel vor.

\begin{DefinitionBox}{Normalenvektor}{analytic_geometry_normal_vector}
  Ein von null verschiedener Vektor $\vec n$ heißt
  \emph{Normalenvektor} \EnglishTerm{normal vector} der Ebene $E$, wenn er zu
  jedem Richtungsvektor der Ebene senkrecht steht.

  Für zwei Richtungsvektoren $\vec a$ und $\vec b$ von $E$ gilt also
  \[
    \vec n\cdot\vec a=0,
    \qquad
    \vec n\cdot\vec b=0.
  \]
\end{DefinitionBox}

\begin{center}
  \CachedTikzFigure{content/resources/chapters/analytic_geometry/tikz/plane_normal_form.tex}
\end{center}

\begin{TheoremBox}{Normalenform einer Ebene}{analytic_geometry_plane_normal_equation}
  Sei $P$ ein Punkt der Ebene $E$ mit Ortsvektor $\vec p$ und sei $\vec n$ ein
  Normalenvektor von $E$. Die \emph{Normalenform}
  \EnglishTerm{normal equation} der Ebene lautet
  \[
    E:\quad (\vec x-\vec p)\cdot\vec n=0.
  \]
  Denn $\vec x-\vec p$ verbindet zwei Punkte der Ebene und steht deshalb
  senkrecht auf $\vec n$.
\end{TheoremBox}

\begin{KeyPointBox}{Koordinatenform und Hessesche Normalform}{analytic_geometry_plane_coordinate_and_hesse_form}
  Für $\vec n=(n_1,n_2,n_3)^{\mathsf T}$ und
  $\vec x=(x_1,x_2,x_3)^{\mathsf T}$ entsteht durch Ausmultiplizieren die
  \emph{Koordinatenform} \EnglishTerm{coordinate equation}
  \[
    E:\quad n_1x_1+n_2x_2+n_3x_3=d,
    \qquad d=\vec n\cdot\vec p.
  \]
  Mit dem Einheitsnormalenvektor
  $\vec e_n=\vec n/\lvert\vec n\rvert$ erhält man die
  \emph{Hessesche Normalform} \EnglishTerm{Hesse normal form}
  \[
    E:\quad (\vec x-\vec p)\cdot\vec e_n=0.
  \]
\end{KeyPointBox}

\begin{RemarkBox}[title=Geraden im zweidimensionalen Raum]
  Im $\R^2$ beschreibt dieselbe Idee eine Gerade:
  \[
    g:\quad (\vec x-\vec p)\cdot\vec n=0.
  \]
  Der Normalenvektor $\vec n$ steht dort senkrecht auf dem Richtungsvektor der
  Geraden.
\end{RemarkBox}
